\documentclass[11pt,a4paper]{article}

\usepackage[margin=2.5cm]{geometry}
\usepackage{amsmath,amssymb,amsthm}
\usepackage{mathtools}
\usepackage{bm}
\usepackage{booktabs}
\usepackage{array}
\usepackage{enumitem}
\usepackage{xcolor}
\usepackage{hyperref}
\usepackage{cleveref}

\hypersetup{
  colorlinks=true,
  linkcolor=blue!50!black,
  urlcolor=blue!50!black,
  citecolor=blue!50!black,
}

\newcommand{\Hop}{\hat{H}}
\newcommand{\Top}{\hat{T}}
\newcommand{\Vop}{\hat{V}}
\newcommand{\Aop}{\hat{\mathcal{A}}}
\newcommand{\Eint}{\Delta E_{\mathrm{int}}}
\newcommand{\Eiso}{E^{(\mathrm{iso})}}
\newcommand{\Ecpl}{E_{\mathrm{complex}}}
\newcommand{\rmax}{r_{\mathrm{max}}}

\title{Energy decomposition of interacting molecular complexes:\\
theory and feasibility within \textsc{MACE}-jax}
\author{}
\date{}

\begin{document}
\maketitle

\begin{abstract}
We examine, starting from the many-body Schr\"odinger equation, the extent to
which the total energy of a complex of interacting molecules can be cleanly
partitioned into self-energies of the individual molecules and inter-molecular
interaction contributions. We give the full many-body expansion (MBE) of the
total energy to all orders, together with its M\"obius-inversion form. We then
ask whether this decomposition is accessible to a trained \textsc{MACE}
equivariant message-passing potential (specifically the present
\texttt{mace-jax} codebase), distinguishing decompositions that follow from a
single forward pass, those that require orchestration across sub-graph forward
passes, and those that would require code changes or are not in principle
accessible to a model of this class.
\end{abstract}

\section{Setting}

Consider a chemical complex composed of $M$ molecules,
$\mathcal{C} = \{1, 2, \dots, M\}$, with nuclei partitioned by molecule,
$\{A\} = \bigsqcup_{\alpha=1}^{M} \mathcal{N}_\alpha$,
and electrons partitioned (only formally; see \cref{sec:obstacles}) as
$\{i\} = \bigsqcup_{\alpha=1}^{M} \mathcal{E}_\alpha$.
We work in the Born--Oppenheimer approximation throughout.

For any non-empty subset of molecules
$S \subseteq \mathcal{C}$, write $E(S)$ for the exact ground-state electronic
energy of the (isolated) sub-system containing the molecules indexed by $S$,
i.e.\ the lowest eigenvalue of the electronic Hamiltonian restricted to nuclei
in $\bigsqcup_{\alpha \in S} \mathcal{N}_\alpha$. We use the abbreviations
$\Ecpl \equiv E(\mathcal{C})$ for the full complex and
$\Eiso_\alpha \equiv E(\{\alpha\})$ for an isolated monomer.

\section{Hamiltonian split}

In atomic units the electronic Hamiltonian of the complex is
\begin{equation}
\Hop \;=\; -\tfrac12 \sum_{i} \nabla_i^2
\;-\; \sum_{i,A} \frac{Z_A}{|\mathbf{r}_i - \mathbf{R}_A|}
\;+\; \sum_{i<j} \frac{1}{|\mathbf{r}_i - \mathbf{r}_j|}
\;+\; \sum_{A<B} \frac{Z_A Z_B}{|\mathbf{R}_A - \mathbf{R}_B|}.
\label{eq:H-full}
\end{equation}
Splitting every two-index sum into ``same-molecule'' and ``different-molecule''
contributions gives
\begin{equation}
\Hop \;=\; \sum_{\alpha=1}^{M} \Hop_\alpha^{(0)} \;+\; \Vop_{\mathrm{int}},
\label{eq:H-split}
\end{equation}
where $\Hop_\alpha^{(0)}$ is the electronic Hamiltonian one would write down for
molecule $\alpha$ in isolation (kinetic energy of its electrons, electron--own-nucleus
attractions, electron--electron repulsion among its electrons, and own
nucleus--nucleus repulsion), and $\Vop_{\mathrm{int}}$ collects all
electron--nucleus, electron--electron and nucleus--nucleus terms in which the
two indices belong to \emph{different} molecules:
\begin{equation}
\Vop_{\mathrm{int}}
\;=\; -\!\!\sum_{\substack{i \in \mathcal{E}_\alpha \\ A \in \mathcal{N}_\beta \\ \alpha \ne \beta}}
        \!\!\frac{Z_A}{|\mathbf{r}_i - \mathbf{R}_A|}
\;+\; \sum_{\substack{i \in \mathcal{E}_\alpha,\, j \in \mathcal{E}_\beta \\ \alpha < \beta}}
        \frac{1}{|\mathbf{r}_i - \mathbf{r}_j|}
\;+\; \sum_{\substack{A \in \mathcal{N}_\alpha,\, B \in \mathcal{N}_\beta \\ \alpha < \beta}}
        \frac{Z_A Z_B}{|\mathbf{R}_A - \mathbf{R}_B|}.
\label{eq:V-int}
\end{equation}
At the operator level, \cref{eq:H-split} is exact and unambiguous.

\section{Obstacles to a clean per-molecule decomposition}
\label{sec:obstacles}

\subsection*{(i) Indistinguishability of electrons}

The exact ground state $\Psi$ of $\Hop$ must be antisymmetric under exchange of
\emph{any} two electrons, including electrons that we would like to assign to
different molecules. There is no quantum-mechanical observable that labels an
electron as belonging to a particular monomer. Two consequences follow:
\begin{itemize}[leftmargin=1.4em]
\item A factorised ansatz $\Psi = \prod_\alpha \Psi_\alpha$ violates Pauli;
\item The antisymmetrised product $\Aop \prod_\alpha \Psi_\alpha$ (Heitler--London
type) introduces exchange/overlap cross-terms that intrinsically mix monomers
and prevent a strict additive decomposition of expectation values.
\end{itemize}
In particular, $\langle \Psi | \Hop_\alpha^{(0)} | \Psi \rangle \neq \Eiso_\alpha$
in general: the operator restricted to molecule $\alpha$ is well defined, but
its expectation value in the supersystem state is not the isolated-monomer
energy.

\subsection*{(ii) Non-uniqueness of further decomposition}

What \emph{is} unambiguous is the difference between the supersystem energy and
the sum of isolated-monomer energies:
\begin{equation}
\Eint \;\equiv\; \Ecpl \;-\; \sum_{\alpha=1}^{M} \Eiso_\alpha.
\label{eq:Eint}
\end{equation}
This quantity is rigorously defined (each piece is a ground-state eigenvalue)
and basis-set-independent in the complete-basis limit. In a finite atom-centred
basis it is contaminated by basis set superposition error (BSSE) and requires
counterpoise correction.

Decomposing $\Eint$ further into physical components
(electrostatic, exchange-repulsion, induction/polarisation, dispersion,
charge-transfer) is \emph{scheme-dependent}:
\begin{itemize}[leftmargin=1.4em]
\item \textbf{SAPT} (symmetry-adapted perturbation theory) treats
$\Vop_{\mathrm{int}}$ perturbatively about the antisymmetrised product of
monomer ground states, yielding an order-by-order series of contributions with
clean physical interpretation.
\item \textbf{Energy decomposition analyses} (Kitaura--Morokuma, ALMO-EDA,
NEDA, \ldots) define intermediate variational states (frozen monomer densities,
polarised but charge-transfer-free orbitals, \ldots) and read off energy
differences. Different EDA schemes generally give numerically different, and
sometimes qualitatively different, component energies for the same complex.
\end{itemize}
What \emph{is} unique, given a single convention, is the decomposition by
cluster size: the many-body expansion.

\section{The many-body expansion (MBE) to all orders}
\label{sec:MBE}

\subsection{Definition by recursion}

Define $n$-body energies by recursion on $|S|$:
\begin{align}
\Delta E^{(1)}_{\{\alpha\}} &\;\equiv\; E(\{\alpha\}) \;=\; \Eiso_\alpha,
\label{eq:DE1} \\[2pt]
\Delta E^{(2)}_{\{\alpha,\beta\}} &\;\equiv\; E(\{\alpha,\beta\})
   \;-\; \Eiso_\alpha \;-\; \Eiso_\beta,
\label{eq:DE2} \\[2pt]
\Delta E^{(3)}_{\{\alpha,\beta,\gamma\}} &\;\equiv\; E(\{\alpha,\beta,\gamma\})
   \;-\!\!\!\sum_{\delta \in \{\alpha,\beta,\gamma\}}\!\!\! \Eiso_\delta
   \;-\!\!\!\!\!\sum_{\substack{T \subset \{\alpha,\beta,\gamma\} \\ |T|=2}}
        \!\!\!\!\!\Delta E^{(2)}_T,
\label{eq:DE3}
\end{align}
and in general, for any non-empty $S \subseteq \mathcal{C}$ with $|S|=n$,
\begin{equation}
\boxed{\;
\Delta E^{(n)}_{S}
\;\equiv\; E(S)
\;-\!\! \sum_{\varnothing \ne T \subsetneq S} \!\! \Delta E^{(|T|)}_{T}.
\;}
\label{eq:DEn-recursive}
\end{equation}
The $n$-body energy $\Delta E^{(n)}_S$ is, by construction, the part of $E(S)$
that is \emph{not} captured by all proper-subset clusters: it is the irreducible
$n$-body correction.

\subsection{Closed form via M\"obius inversion}

Equation \cref{eq:DEn-recursive} is the M\"obius inversion of $E(S)$ with
respect to set inclusion. Solving it gives the explicit closed form
\begin{equation}
\boxed{\;
\Delta E^{(n)}_{S}
\;=\; \sum_{T \subseteq S} (-1)^{|S|-|T|} \, E(T),
\quad |S|=n.
\;}
\label{eq:DEn-mobius}
\end{equation}
Conversely, the energy of any sub-cluster is recovered by summing $n$-body
contributions over its own subsets,
\begin{equation}
E(S) \;=\; \sum_{\varnothing \ne T \subseteq S} \Delta E^{(|T|)}_{T}.
\label{eq:E-from-DE}
\end{equation}

\subsection{Total-energy expansion}

Specialising \cref{eq:E-from-DE} to $S = \mathcal{C}$ gives the many-body
expansion of the full complex,
\begin{equation}
\boxed{\;
\Ecpl
\;=\; \sum_{\alpha} \Eiso_\alpha
\;+\; \sum_{\alpha<\beta} \Delta E^{(2)}_{\{\alpha,\beta\}}
\;+\; \sum_{\alpha<\beta<\gamma} \Delta E^{(3)}_{\{\alpha,\beta,\gamma\}}
\;+\; \cdots
\;+\; \Delta E^{(M)}_{\mathcal{C}}.
\;}
\label{eq:MBE-total}
\end{equation}
Equivalently, separating self-energies from interactions,
\begin{equation}
\Eint
\;=\; \Ecpl - \sum_{\alpha} \Eiso_\alpha
\;=\; \sum_{n=2}^{M} \;\sum_{\substack{S \subseteq \mathcal{C} \\ |S|=n}}
        \Delta E^{(n)}_{S}.
\label{eq:Eint-MBE}
\end{equation}

\subsection{Properties}

\begin{itemize}[leftmargin=1.4em]
\item \emph{Exactness.} The expansion in \cref{eq:MBE-total} is identically
exact at any order, provided \emph{all} terms are retained: it terminates at
$n = M$.
\item \emph{Permutation symmetry.} Each $\Delta E^{(n)}_{S}$ depends only on the
set $S$, not on a labelling of its elements; \cref{eq:DEn-mobius} is manifestly
symmetric.
\item \emph{Counting.} Truncating at order $n^\star$ requires
$\sum_{n=1}^{n^\star} \binom{M}{n}$ sub-system calculations:
$M$ monomers, $\binom{M}{2}$ dimers, $\binom{M}{3}$ trimers, etc.
\item \emph{Convergence.} For systems dominated by short-ranged, weakly
polarising interactions (e.g.\ hydrogen-bonded clusters, van der Waals
complexes), the series converges quickly, with $\Delta E^{(2)}$ accounting for
the bulk, $\Delta E^{(3)}$ giving cooperative/anti-cooperative corrections at
the few-percent level, and $\Delta E^{(n \ge 4)}$ usually negligible. For
metallic, highly polarisable, or strongly delocalised systems convergence is
slow and truncation can be unreliable.
\item \emph{BSSE.} If $E(S)$ is computed in a finite atom-centred basis, each
sub-system uses a different basis set and \cref{eq:DEn-mobius} mixes
inconsistent bases. Counterpoise correction (computing every $E(T)$ in the
basis of the largest cluster of interest) restores consistency at extra cost.
\end{itemize}

\subsection{Worked example: $M = 4$}

For a four-molecule complex $\mathcal{C} = \{1,2,3,4\}$, \cref{eq:MBE-total}
reads
\begin{equation}
\begin{aligned}
\Ecpl \;=\;& \sum_{\alpha=1}^{4} \Eiso_\alpha
\;+\; \sum_{1 \le \alpha < \beta \le 4} \Delta E^{(2)}_{\{\alpha,\beta\}}
\;+\; \sum_{1 \le \alpha < \beta < \gamma \le 4} \Delta E^{(3)}_{\{\alpha,\beta,\gamma\}}
\;+\; \Delta E^{(4)}_{\{1,2,3,4\}},
\end{aligned}
\end{equation}
with, for instance,
\begin{align}
\Delta E^{(2)}_{\{1,2\}} \;=\;& E(\{1,2\}) - \Eiso_1 - \Eiso_2, \\
\Delta E^{(3)}_{\{1,2,3\}} \;=\;& E(\{1,2,3\})
   - \!\!\sum_{\alpha \in \{1,2,3\}} \!\! \Eiso_\alpha
   - \!\!\!\!\sum_{\substack{T \subset \{1,2,3\} \\ |T|=2}} \!\!\!\! \Delta E^{(2)}_T, \\
\Delta E^{(4)}_{\{1,2,3,4\}} \;=\;& \sum_{T \subseteq \{1,2,3,4\}} (-1)^{4-|T|} E(T) \nonumber\\
   \;=\;& E(\{1,2,3,4\}) \nonumber\\
   &\;-\; \big[E(\{1,2,3\}) + E(\{1,2,4\}) + E(\{1,3,4\}) + E(\{2,3,4\})\big] \nonumber\\
   &\;+\; \big[E(\{1,2\}) + E(\{1,3\}) + E(\{1,4\}) + E(\{2,3\}) + E(\{2,4\}) + E(\{3,4\})\big] \nonumber\\
   &\;-\; \big[\Eiso_1 + \Eiso_2 + \Eiso_3 + \Eiso_4\big].
\end{align}
The pattern of alternating signs and binomial coefficients in
\cref{eq:DEn-mobius} is visible explicitly: $\binom{4}{0}{=}\;0$ empty term,
$\binom{4}{1}=4$ monomer terms with sign $-$, $\binom{4}{2}=6$ dimer terms with
sign $+$, $\binom{4}{3}=4$ trimer terms with sign $-$, and the supersystem with
sign $+$.

\subsection{Summary of what is and is not unique}

\begin{table}[h]
\centering
\renewcommand{\arraystretch}{1.15}
\begin{tabular}{lll}
\toprule
Quantity & Status & Caveat \\
\midrule
$\Ecpl$ and $\Eiso_\alpha$            & QM observable               & complete basis \\
$\Eint = \Ecpl - \sum_\alpha \Eiso_\alpha$ & rigorous, basis-set-dep.\ at finite basis & BSSE / counterpoise \\
$\Delta E^{(n)}_{S}$ for $n \ge 2$    & rigorous (M\"obius)         & idem. \\
electrostatic / exchange / dispersion / \ldots & scheme-dependent (SAPT, EDA) & no canonical choice \\
``energy of molecule $\alpha$ inside $\mathcal{C}$'' & not a QM observable    & only differences with $\Eiso_\alpha$ are \\
\bottomrule
\end{tabular}
\end{table}

\section{Feasibility within \textsc{MACE}-jax}

\textsc{MACE} is a strictly local, equivariant message-passing graph neural
network potential. The trained model $f_\theta$ takes a single graph
$G$ (atom positions, species, periodic cell, edges within a fixed cutoff
$\rmax$) and returns a scalar energy $E_\theta(G)$ together with per-atom
contributions $\varepsilon_i$ that satisfy
$E_\theta(G) = \sum_{i \in G} \varepsilon_i$ by construction.

\subsection{Architectural facts relevant to decomposition}

The following points are observed in the present codebase
(\texttt{mace\_jax/modules/models.py}, \texttt{mace\_jax/modules/blocks.py},
\texttt{mace\_jax/data/neighborhood.py}, \texttt{mace\_jax/data/utils.py}):

\begin{enumerate}[leftmargin=1.6em]
\item \emph{Per-atom output, summed to total.} $L$ message-passing layers each
emit a per-atom scalar via a readout block; these are stacked and summed to
yield $\varepsilon_i$, and atoms are scatter-summed by graph index to yield the
graph-level total.

\item \emph{Atomic reference energies.} A per-species offset $E_0^{Z_i}$ is
subtracted, with $\mathrm{stop\_gradient}$ applied so it is not learnt during
training. The model already exposes
\[
E_\theta^{\mathrm{int}}(G) \;=\; E_\theta(G) - \sum_{i \in G} E_0^{Z_i},
\]
but here ``interaction'' means total minus per-atom self-energy; it is
\emph{not} an inter-molecular interaction energy in the sense of
\cref{eq:Eint}.

\item \emph{Strict locality.} Each message-passing layer aggregates only over
neighbours within $\rmax$. With $L$ layers, the energy of atom $i$ depends only
on atoms within graph distance $L$, hence within Euclidean distance
$L \cdot \rmax$ at most. The default is $L = 3$.

\item \emph{No fragment indexing.} Inputs carry a \texttt{batch} array that
groups atoms into \emph{independent} structures, but no per-atom
``molecule'' or ``fragment'' index within a structure.
\end{enumerate}

\subsection{What the model can already compute}

The supersystem decomposition \cref{eq:Eint} and the full MBE \cref{eq:MBE-total}
are accessible without any code changes, provided one is willing to perform one
forward pass per cluster:

\begin{itemize}[leftmargin=1.4em]
\item Build the graph for each subset $S \subseteq \mathcal{C}$ of interest
(using only the atoms in $\bigsqcup_{\alpha \in S} \mathcal{N}_\alpha$;
the neighbour-list builder is given the corresponding positions and the same
$\rmax$).
\item Evaluate $\hat E_\theta(S) \equiv f_\theta(\text{graph}(S))$, treating
each $\hat E_\theta(S)$ as the model's surrogate for $E(S)$.
\item Combine via \cref{eq:DEn-mobius,eq:MBE-total} to obtain
$\widehat{\Delta E}^{(n)}_{S}$ and $\widehat{\Eint}$.
\end{itemize}
All sub-graphs at a given truncation order can be packed into a single padded
batch, so practical cost is $\mathcal{O}(\sum_{n \le n^\star} \binom{M}{n})$
forward passes.

\subsection{What requires modest extra code}

\paragraph{Model-internal partition by fragment.}
Because \textsc{MACE} produces per-atom energies $\varepsilon_i$, one can also
define, from a \emph{single} forward pass on the complex,
\begin{equation}
\widetilde{E}_\alpha \;\equiv\; \sum_{i \in \mathcal{N}_\alpha} \varepsilon_i\big(\mathrm{graph}(\mathcal{C})\big),
\qquad
\sum_{\alpha} \widetilde{E}_\alpha = \Ecpl.
\label{eq:tilde-E}
\end{equation}
This implements an EDA-style ``share'' of the supersystem energy assigned to
each molecule \emph{in situ}. It is not a quantum-mechanical observable, and it
depends on architectural and training choices; two \textsc{MACE} models trained
to identical totals will in general produce different per-atom partitions.
Implementation needs:
\begin{itemize}[leftmargin=1.4em]
\item add an optional \texttt{fragment\_index} field to the data dict
(length equal to \texttt{num\_atoms}); construct it during graph building;
\item in \texttt{models.py}, after the existing per-graph
\texttt{scatter\_sum}, perform a second \texttt{scatter\_sum} of the per-atom
energies grouped by \texttt{fragment\_index} to produce
$\widetilde{E}_\alpha$ as an additional output.
\end{itemize}
This is on the order of 10--20 lines and does not alter training.

\paragraph{Many-body driver.}
A small driver utility that, given a complex with a fragment partition,
enumerates subsets up to a chosen $n^\star$, builds their graphs, batches the
forward passes, and applies \cref{eq:DEn-mobius} to return
$\{\widehat{\Delta E}^{(n)}_S\}$. Pure orchestration, no model changes.

\subsection{What is not in principle accessible}

\begin{itemize}[leftmargin=1.4em]
\item \emph{Component decomposition} of $\Eint$ into electrostatic, exchange,
induction, dispersion, and charge-transfer in the SAPT/EDA sense. These rely
on intermediate, scheme-defined wavefunctions that the model does not
construct. \textsc{MACE} returns a single learnt scalar; no architectural
modification short of a different model class can recover those components
post hoc.

\item \emph{Counterpoise correction} of $\widehat{\Eint}$. \textsc{MACE} does
not use atom-centred Gaussian basis sets, so traditional BSSE does not arise
at inference time. However, the model has been trained against reference
energies (DFT, CCSD(T), \ldots) that may themselves be BSSE-contaminated;
this is a property of the training data and cannot be repaired by inference-time
manipulation.

\item \emph{Long-range physics beyond the receptive field.} Two molecules
separated by more than $L \cdot \rmax$ are non-adjacent in the message-passing
graph, and the model's predicted $\widehat{\Eint}$ between them is identically
zero up to the per-species offsets. Long-range electrostatics and dispersion
are captured only insofar as the model has been augmented with explicit
long-range heads (learnt charges + Ewald summation, learnt dispersion
coefficients, \ldots). Whether this is the case is a property of the model
variant and the training set, and should be checked against the typical
inter-molecular distances in the complex of interest before assigning physical
meaning to a computed $\widehat{\Eint}$.
\end{itemize}

\subsection{Recommended workflow}

For a complex $\mathcal{C}$ with $M$ molecules, given a trained
\textsc{MACE}-jax model:
\begin{enumerate}[leftmargin=1.6em]
\item Verify that the largest inter-molecular distance of interest does not
exceed $L \cdot \rmax$, or that an appropriate long-range correction is
present.
\item Compute $\widehat{E}_\theta(\mathcal{C})$ on the full complex graph.
\item Compute $\widehat{\Eiso}_\alpha = \widehat{E}_\theta(\{\alpha\})$ for
each monomer; report
$\widehat{\Eint} = \widehat{E}_\theta(\mathcal{C}) - \sum_\alpha \widehat{\Eiso}_\alpha$.
\item For a finer view, evaluate dimers and (if needed) trimers, and report
$\widehat{\Delta E}^{(2)}_{\{\alpha,\beta\}}$ and
$\widehat{\Delta E}^{(3)}_{\{\alpha,\beta,\gamma\}}$ via
\cref{eq:DEn-mobius}; the residual
$\widehat{\Eint} - \sum_{|S|=2}\widehat{\Delta E}^{(2)}_S - \sum_{|S|=3}\widehat{\Delta E}^{(3)}_S$
diagnoses higher-body cooperativity within the model.
\item Optionally, for an in-situ ``share'' per molecule, add the
\texttt{fragment\_index} hook described above and report $\widetilde{E}_\alpha$
from \cref{eq:tilde-E}, with the explicit caveat that it is a model-internal
partition rather than a quantum observable.
\end{enumerate}

\section{Effective surfaces from a focus set}
\label{sec:effective}

\subsection{Motivation}

When studying the conformational landscape of a complex it is often the case
that not all minima are physically interesting. A common situation is that a
designated molecule, say molecule $1$, is the actor of interest and one wishes
to scan the energy as a function of the poses of $\{2,\dots,M\}$ around it.
Configurations in which two spectator molecules cluster strongly with each
other but stay far from molecule $1$ produce deep minima in $\Ecpl$ that are
irrelevant to the question being asked.

The MBE machinery of \cref{sec:MBE} suggests an obvious filter: keep only the
many-body terms whose subset $S$ contains molecule $1$, and discard those
confined to spectators. More generally, given a \emph{focus set}
$F \subseteq \mathcal{C}$, define
\begin{equation}
E_{\mathrm{eff}}^{F}
\;\equiv\; \!\!\!\!\sum_{\substack{S \subseteq \mathcal{C}\\ S \cap F \ne \varnothing}}
        \!\!\!\! \Delta E^{(|S|)}_{S}.
\label{eq:Eeff-MBE}
\end{equation}
Every term that survives in $E_{\mathrm{eff}}^{F}$ involves at least one
fragment of interest; every spectator-only cluster has been cancelled out.

\subsection{Closed form}

The same identity \cref{eq:E-from-DE} that produces the MBE also collapses
\cref{eq:Eeff-MBE}. Apply it once to $S = \mathcal{C}$ and once to
$S = \mathcal{C} \setminus F$ and subtract:
\begin{equation}
\boxed{\;
E_{\mathrm{eff}}^{F}
\;=\; E(\mathcal{C}) \;-\; E(\mathcal{C} \setminus F),
\;}
\label{eq:Eeff-closed}
\end{equation}
because the cluster sums on the right share exactly the spectator-confined
terms $\Delta E^{(|T|)}_T$ with $T \cap F = \varnothing$, which cancel pairwise.

If one further wants the surface to vanish in the non-interacting limit
(rather than approaching $\sum_{\alpha \in F} \Eiso_\alpha$), subtract the
focus-only energy:
\begin{equation}
E_{\mathrm{eff}}^{F,0}
\;\equiv\; E(\mathcal{C}) \;-\; E(\mathcal{C} \setminus F) \;-\; E(F)
\;=\; \!\!\!\!\sum_{\substack{S \subseteq \mathcal{C}\\ S \cap F \ne \varnothing\\ S \not\subseteq F}}
       \!\!\!\! \Delta E^{(|S|)}_{S}.
\label{eq:Eeff0-closed}
\end{equation}

\subsection{Properties}

\begin{itemize}[leftmargin=1.4em]
\item \emph{Spectator gauge invariance.}\, Translating, rotating, or rearranging
the spectators \emph{among themselves} while keeping each of them outside the
receptive field of every fragment in $F$ leaves $E_{\mathrm{eff}}^{F,0}$
identically zero (in $E_{\mathrm{eff}}^{F}$, the constant
$\sum_{\alpha \in F}\Eiso_\alpha$ is unaffected). Spectator clustering minima
do not appear in the surface.

\item \emph{Cooperativity is preserved.}\, Three-body and higher contributions
$\Delta E^{(n)}_{S}$ with $S \cap F \ne \varnothing$ remain in the sum at every
order. The construction is not a 2-body truncation: many-body cooperativity
mediated through fragments in $F$ is captured exactly.

\item \emph{Cost.}\, Two energy evaluations
($E(\mathcal{C})$, $E(\mathcal{C}\setminus F)$) for $E_{\mathrm{eff}}^{F}$, or
three for $E_{\mathrm{eff}}^{F,0}$. No subset enumeration; cost is
$\mathcal{O}(1)$ in $M$ at fixed $|F|$.

\item \emph{Forces.}\, Applying $\nabla_{\mathbf r_i}$ to
\cref{eq:Eeff-closed} gives, for $i \in \mathcal{N}_\alpha$ with
$\alpha \notin F$,
\[
-\nabla_{\mathbf r_i} E_{\mathrm{eff}}^{F}
\;=\; \mathbf F^{\mathcal{C}}_i \;-\; \mathbf F^{\mathcal{C}\setminus F}_i,
\]
i.e.\ the force on a spectator atom is the in-complex force minus the
spectator-only force, leaving only the part mediated by $F$. For $i$ in a
focus fragment, $E(\mathcal{C}\setminus F)$ does not depend on $\mathbf r_i$
and the spectator-only term contributes nothing.
\end{itemize}

\subsection{Behaviour outside the receptive field}

Two atoms separated by graph distance greater than $L$ in the message-passing
graph do not communicate. If every atom of $F$ is more than $L \cdot \rmax$
away from every atom of a particular spectator fragment $\beta$, then in
$E(\mathcal{C})$ the contributions involving atoms of $\beta$ depend only on
$\beta$'s own neighbourhood (which is a subset of $\mathcal{C}\setminus F$),
and they are matched exactly by the corresponding contributions in
$E(\mathcal{C}\setminus F)$. Hence the difference \cref{eq:Eeff-closed} is
flat in any direction that moves $\beta$ around without bringing it within
$L \cdot \rmax$ of $F$. This is the desired behaviour. In particular,
neighbour-list rebuilds in $\mathcal{C}\setminus F$ produce no spurious
edges --- edges that crossed the focus/spectator boundary in $\mathcal{C}$
are simply absent in $\mathcal{C}\setminus F$.

\subsection{Implementation: \texttt{mace\_jax/effective.py}}

The construction has been wrapped as \texttt{mace\_jax.effective.EffectiveEnergy}.
Given a fitted model and parameters with the same signature used by the
existing ASE calculator,
\begin{verbatim}
from mace_jax.effective import EffectiveEnergy

eff = EffectiveEnergy(model=model_fn, params=params, r_max=r_max,
                      atomic_numbers=z_list)

# fragment_index: per-atom integer fragment id (length = len(atoms)).
# focus: iterable of fragment ids forming F.
result = eff(atoms, fragment_index, focus={0}, baseline_subtract=True)

result.energy        # E_eff^F (or E_eff^{F,0})
result.forces        # shape (N, 3), in the original atoms layout
result.components    # {'E_full': ..., 'E_complement': ..., 'E_focus': ...}
\end{verbatim}
Internally the module builds the two (or three) sub-system graphs by
re-running the existing neighbour-list builder on the appropriate atom
subsets, jit-compiles the predictor once, pads the graphs with a per-shape
edge budget (so nearby geometries reuse the same compiled trace), and
combines the energies and forces according to \cref{eq:Eeff-closed} or
\cref{eq:Eeff0-closed}. Forces from each sub-system are scattered back into
the full-atoms layout before being subtracted, so the returned array is
indexed exactly as the input \texttt{atoms} object.

\section{Conclusion}

The two-way split of the total energy of a molecular complex into a sum of
isolated-monomer energies and an interaction energy is theoretically
unambiguous. Its refinement into a many-body expansion is also unique once the
recursive convention \cref{eq:DEn-recursive} (equivalently, the M\"obius form
\cref{eq:DEn-mobius}) is fixed, and is exact at every truncation order if all
terms up to that order are retained. Further decomposition of the interaction
energy into physically named components is not unique and requires choosing a
SAPT/EDA scheme.

For \textsc{MACE}-jax, the supersystem decomposition \cref{eq:Eint} and the
full many-body expansion \cref{eq:MBE-total} are computable using only the
existing forward pass, with no architectural changes; only orchestration over
sub-graphs is needed. A model-internal per-fragment partition
\cref{eq:tilde-E} is straightforward to expose with a small change to the data
dict and the energy aggregation step. Component-style decompositions
(electrostatic / exchange / dispersion / \ldots) are not accessible from a
single learnt scalar potential and would require a different model class. The
fidelity of any such decomposition is bounded by the receptive field
$L \cdot \rmax$ and by whether the underlying training data and model variant
account for long-range interactions.

\end{document}
